\section{Introduction}

The transition to sustainable architecture and energy-efficient building operations is increasingly recognized as a critical global priority. Commercial buildings account for a substantial portion of global energy consumption, and lighting alone typically represents between 15\% and 20\% of a commercial facility's total electricity use \cite{eia_lighting}. In response, modern building codes and green certification programs (e.g., LEED, BREEAM) have mandated stringent energy performance standards during the design phase. However, a pervasive challenge remains: the ``verification gap'' or ``performance gap'' \cite{performance_gap_1}. This gap refers to the well-documented discrepancy between the highly optimized energy efficiency predicted during a building's design phase and the often significantly higher actual energy consumption observed once the building is operational \cite{performance_gap_2}.

A primary driver of this verification gap is suboptimal operational control, particularly in systems like lighting and HVAC (Heating, Ventilation, and Air Conditioning), which are highly sensitive to occupant behavior \cite{performance_gap_3}. While modern facilities are equipped with sophisticated Building Management Systems (BMS), these systems frequently rely on static schedules that fail to adapt to dynamic occupancy patterns, leading to substantial energy waste---such as fully illuminating empty floors during evenings or weekends.

Addressing this inefficiency requires granular, continuous auditing of specific end-use loads. Unfortunately, direct measurement via hardware sub-metering (installing dedicated physical meters on every lighting circuit) is prohibitively expensive, disruptive, and logistically complex for the vast majority of existing commercial building stock \cite{nilm_survey_1}. Consequently, facility managers are left with whole-building aggregate smart meter data, which provides a macro-level view of consumption but obscures the specific operational patterns of individual systems like lighting.

To bridge this gap, Non-Intrusive Load Monitoring (NILM)---also known as energy disaggregation---has emerged as a powerful computational alternative to hardware sub-metering \cite{hart_1992}. NILM frames the extraction of individual appliance loads from an aggregate power signal as a blind source separation problem. While high-frequency NILM (sampling at kHz to MHz rates) has shown success in identifying transient electrical signatures of specific devices \cite{nilm_survey_2}, commercial smart meters typically report data at much lower frequencies, such as hourly intervals. At this resolution, transient electrical signatures are lost, and disaggregation must rely entirely on macroscopic temporal patterns.

Recent advancements in deep learning have significantly advanced the state-of-the-art in low-frequency NILM. Architectures such as Convolutional Neural Networks (CNNs) \cite{zhang_2018}, Recurrent Neural Networks (RNNs), and Transformers \cite{bert4nilm} have demonstrated the ability to extract complex temporal dependencies from aggregate load profiles. However, many of these approaches are heavily supervised, requiring large datasets of paired aggregate and sub-metered ground-truth data for training. In the context of commercial buildings, such ground-truth data is exceedingly rare, rendering supervised approaches impractical for widespread deployment.

This paper proposes a fully unsupervised, two-stage deep learning framework designed to audit lighting efficiency in commercial buildings using only standard, hourly aggregate electricity meter data. Our approach does not require any historical sub-metered data or hardware retrofits. Instead, it leverages the inherent structural differences between continuous base loads (e.g., HVAC, servers) and discrete, occupancy-driven loads (e.g., lighting) to perform the disaggregation.

The core contributions of this work are as follows:
\begin{enumerate}
    \item We introduce an unsupervised feature learning pipeline utilizing a 1D Convolutional Autoencoder (1D-CAE) to project daily aggregate load profiles into a compact latent space. This enables robust, scalable clustering of building operations into distinct control regimes (e.g., scheduled weekday use versus reactive weekend use).
    \item We adapt the N-BEATS (Neural Basis Expansion Analysis for Interpretable Time Series) architecture \cite{nbeats_2020}---originally designed for forecasting---as a powerful unsupervised signal decomposition tool. Specialized trend and seasonality blocks isolate high-frequency lighting pulses from the low-frequency HVAC and base load components.
    \item We rigorously validate our framework on the public ASHRAE Great Energy Predictor III dataset \cite{ashrae_2020}, processing nearly 500,000 daily load profiles from 1,448 buildings across 16 climate zones. Our evaluation includes comprehensive baselines, a cross-climate generalization test, and a sensitivity analysis demonstrating the model's robustness to input noise.
    \item We provide a practical, end-to-end case study on an educational facility, illustrating how the extracted lighting profiles can be directly used by facility managers to identify actionable energy-saving opportunities, thereby directly addressing the verification gap.
\end{enumerate}

The remainder of this paper is structured as follows. Section 2 reviews related work in building energy performance, NILM, and deep learning for time series. Section 3 details the proposed two-stage methodology. Section 4 presents the experimental setup, baseline comparisons, and the practical case study. Finally, Section 5 concludes the paper and discusses avenues for future research.

To ensure transparency and reproducibility---a critical requirement for the applied data science community---all code, trained model weights, and an interactive demonstration application are made publicly available at \url{https://github.com/AnonAuthor/FTC-NILM-Auditing}.
